\chapter{Implementazione e risultati degli esperimenti}%
\label{cha:Implementazione e risultati degli esperimenti}
In questo capitolo vengono illustrate le scelte implementative adottate,
il metodo impiegato per eseguire i test e i risultati ottenuti per le varie
configurazioni di parametri.

\section{Configurazione e metodologia dei test}%
\label{sec:Configurazione e metodologia dei test}

\subsection{Insieme delle istanze considerate}%
\label{sub:Insieme delle istanze considerate}
Le performance dell'algoritmo sono state ricavate andandolo a testare
sull'insieme di problemi fornito dalla libreria MIPLIB 2017 \cite{miplib2017}.
In particolare sono stati utilizzati i problemi appartenenti al \emph{Benchmark Set}
nel quale si trovano:
\begin{itemize}
    \item 220 istanze facili;
    \item 19 istanze difficili;
    \item una istanza "aperta".
\end{itemize}
La difficoltà di un'istanza è definita da MIPLIB in base al tempo necessario
per raggiungere l'ottimalità.
Le istanze classificate come "aperte" non sono state ancora state risolte.

\subsection{Software e hardware utilizzato}%
\label{sub:Software utilizzato}
L'algoritmo è stato realizzato utilizzando il linguaggio di programmazione C.
Per la manipolazione e la risoluzione di problemi MIP e LP sono state
utilizzate le funzioni offerte dalla \emph{CPLEX Callable Library} 22.1.1
\cite{cplexCallableLib}.
I test sono stati effettuati sul cluster del dipartimento di ricerca operativa
del Dipartimento di Ingegneria dell'Informazione.
La macchina utilizzata dispone di un processore Intel(R) Xeon(R) CPU E3-1220
V2 @ 3.10GHz e 16GB di RAM.

\subsection{Procedimento seguito e statistiche osservate}%
\label{sub:Procedimento seguito e statistiche osservate}
Trattandosi di un algoritmo euristico non è possibile dimostrare rigorosamente
che una combinazione di valori assegnati ai parametri è migliore rispetto ad
un altra.
Si tratta quindi di verificare in via sperimentale quale combinazione
fornisce i risultati migliori.
Le statistiche che sono state osservate sono le seguenti:
\begin{itemize}
    \item Numero di problemi per cui è stata trovata una soluzione ammissibile per MIP.
    \item Numero di iterazioni medie per ottenere una soluzione ammissibile per MIP.
    \item Tempo medio per ottenere una soluzione ammissibile per MIP.
    \item Qualità media della soluzione trovata.
\end{itemize}
Inizialmente i parametri sono stati fissati a dei valori arbitrari,
successivamente sono state introdotte delle modifiche con l'intento di
migliorare uno o più aspetti di quelli sopra elencati.

Data una soluzione $x$ per un problema MIP avente soluzione ottima $\hat x$,
la qualità $\gamma(x)\in [0,\,1]$ di $x$ è definita come segue:
\begin{equation}
    \gamma(x) =
    \begin{cases}
        0 & \text{se } |c^t\hat x| = |c^t x| = 0\\
        1 & \text{se } c^t \hat x \cdot c^t x < 0\\
        \frac{|c^t \hat x - c^t x|}{\max\{|c^t \hat x|,\, |c^t x|\}} & \text{altrimenti} 
    \end{cases}
\end{equation}
